% =============================================================================
% TCC — ContAR: Plataforma No-Code para Autoria de Narradores Virtuais 3D
%        em Realidade Aumentada Educacional
%
% Autor  : Jorge Luiz Cunha de Freitas
% Curso  : [Nome do Curso] — Centro de Informática (CIn) / UFPE
% Ano    : 2026
% =============================================================================

\documentclass[
  12pt,
  openright,
  twoside,
  a4paper,
  english,
  brazil
]{abntex2}

% ---------------------------------------------------------------------------
% Pacotes
% ---------------------------------------------------------------------------
\usepackage[T1]{fontenc}
\usepackage[utf8]{inputenc}
\usepackage{lmodern}
\usepackage{microtype}
\usepackage{graphicx}
\usepackage{booktabs}
\usepackage{tabularx}
\usepackage{multirow}
\usepackage{xcolor}
\usepackage{listings}
\usepackage{hyperref}
\usepackage{url}
\usepackage{amsmath}
\usepackage{float}
\usepackage{indentfirst}

% ---------------------------------------------------------------------------
% Configurações de código-fonte (para trechos de código)
% ---------------------------------------------------------------------------
\lstset{
  basicstyle=\footnotesize\ttfamily,
  breaklines=true,
  frame=single,
  numbers=left,
  numberstyle=\tiny,
  keywordstyle=\color{blue},
  commentstyle=\color{gray},
  stringstyle=\color{red!60!black},
  backgroundcolor=\color{gray!5},
}

% ---------------------------------------------------------------------------
% Informações do documento
% ---------------------------------------------------------------------------
\titulo{ContAR: Uma Plataforma No-Code para Autoria de Narradores Virtuais 3D
        em Realidade Aumentada Educacional}
\autor{Jorge Luiz Cunha de Freitas}
\local{Recife}
\data{2026}
\orientador{Prof. Dr. [Nome do Orientador]}
% \coorientador{Prof. Dr. [Nome do Coorientador]}  % descomente se houver
\instituicao{
  Universidade Federal de Pernambuco \\
  Centro de Informática (CIn)
}
\tipotrabalho{Trabalho de Conclusão de Curso}
\preambulo{
  Trabalho de Conclusão de Curso apresentado ao Centro de Informática
  da Universidade Federal de Pernambuco como requisito parcial para
  obtenção do grau de [Bacharel em Ciência da Computação].
}

% ---------------------------------------------------------------------------
% Hyperref
% ---------------------------------------------------------------------------
\hypersetup{
  pdftitle={\imprimirtitulo},
  pdfauthor={\imprimirautor},
  colorlinks=true,
  linkcolor=black,
  citecolor=black,
  urlcolor=blue
}

% =============================================================================
\begin{document}
% =============================================================================

\frenchspacing

% ---------------------------------------------------------------------------
% Elementos pré-textuais
% ---------------------------------------------------------------------------
\pretextual

% Capa
\imprimircapa

% Folha de rosto
\imprimirfolhaderosto*

% Ficha catalográfica — gerada eletronicamente pela UFPE
% (substituir pelo PDF gerado no portal da biblioteca)
\begin{fichacatalografica}
  \vspace*{\fill}
  \begin{center}
    \fbox{
      \begin{minipage}[c][8cm]{13cm}
        \small
        Freitas, Jorge Luiz Cunha de.\\
        \hspace*{0.5cm}
        ContAR: Uma Plataforma No-Code para Autoria de Narradores Virtuais 3D
        em Realidade Aumentada Educacional /
        Jorge Luiz Cunha de Freitas. --- Recife, 2026.\\
        \hspace*{0.5cm}
        [Total de páginas] p.\\[0.5em]
        \hspace*{0.5cm}
        Orientador: [Nome do Orientador].\\
        Trabalho de Conclusão de Curso (Graduação) ---
        Universidade Federal de Pernambuco,
        Centro de Informática, [Curso], 2026.\\[0.5em]
        \hspace*{0.5cm}
        1.~Realidade Aumentada.
        2.~WebAR.
        3.~Avatares 3D.
        4.~Síntese de Voz.
        5.~No-Code.
        6.~Educação.
        I.~[Nome do Orientador], orient.
        II.~Título.
      \end{minipage}
    }
  \end{center}
\end{fichacatalografica}

% Folha de aprovação
\begin{folhadeaprovacao}
  \begin{center}
    {\ABNTEXchapterfont\large\imprimirautor}
    \vspace*{\fill}\vspace*{\fill}
    \begin{center}
      \ABNTEXchapterfont\bfseries\Large\imprimirtitulo
    \end{center}
    \vspace*{\fill}
    \hspace{.45\textwidth}
    \begin{minipage}{.5\textwidth}
      \imprimirpreambulo
    \end{minipage}%
    \vspace*{\fill}
  \end{center}

  Aprovado em: \underline{\hspace{3cm}} de \underline{\hspace{3cm}} de 2026.

  \vspace{1cm}
  \noindent\textbf{Banca Examinadora:}

  \assinatura{\textbf{\imprimirorientador} \\ Orientador}
  \assinatura{\textbf{[Membro da Banca 2]} \\ [Instituição]}
  \assinatura{\textbf{[Membro da Banca 3]} \\ [Instituição]}
\end{folhadeaprovacao}

% Dedicatória (opcional)
\begin{dedicatoria}
  \vspace*{\fill}
  \centering
  \noindent
  \textit{A todos os professores e educadores que acreditam\\
  que a tecnologia deve servir às pessoas,\\
  e não o contrário.}
  \vspace*{\fill}
\end{dedicatoria}

% Agradecimentos (opcional)
\begin{agradecimentos}
  Agradeço ao meu orientador, Prof. Dr. [Nome], pelo acompanhamento, pelas
  discussões que moldaram as decisões arquiteturais desta pesquisa e pela
  sugestão de explorar a fronteira entre Inteligência Artificial e narrativas
  educacionais imersivas.

  Aos professores e especialistas que participaram do Expert Walkthrough,
  cujas percepções foram fundamentais para a avaliação crítica da plataforma.

  À comunidade open-source — especialmente os mantenedores do Three.js,
  \texttt{@pixiv/three-vrm} e AR.js — cujas contribuições tornaram este
  trabalho tecnicamente viável.

  À UFPE e ao CIn, pelo ambiente de pesquisa e pela infraestrutura que
  possibilitaram o desenvolvimento deste projeto.

  E à minha família, pelo apoio incondicional durante toda a graduação.
\end{agradecimentos}

% Epígrafe (opcional)
\begin{epigrafe}
  \vspace*{\fill}
  \begin{flushright}
    \textit{``Any sufficiently advanced technology\\
    is indistinguishable from magic.''}\\
    --- Arthur C. Clarke
  \end{flushright}
\end{epigrafe}

% Resumo em português
\begin{resumo}
  A Realidade Aumentada (RA) tem sido amplamente investigada como recurso
  pedagógico capaz de promover experiências de aprendizagem interativas e
  multimodais. No entanto, a criação de conteúdo personalizado em RA ainda
  impõe barreiras técnicas significativas a educadores sem formação em
  desenvolvimento de software ou modelagem tridimensional. Um estudo
  exploratório anterior, realizado como parte desta pesquisa, identificou
  que ferramentas disponíveis no mercado operam de forma fragmentada,
  exigindo que o usuário domine múltiplas plataformas independentes para
  concluir um fluxo completo de autoria em RA.

  Diante dessa lacuna, este trabalho apresenta o \textbf{ContAR}, uma
  plataforma web no-code para autoria de narradores virtuais 3D em Realidade
  Aumentada educacional. A plataforma integra, em uma única interface, a
  criação e customização de avatares 3D, a geração de voz sintética com
  sincronização labial precisa via visemas, a composição de cenas e histórias
  sequenciais, e a publicação em RA acessível diretamente pelo navegador
  (WebXR Surface AR e AR.js Marker AR).

  A arquitetura adota um modelo SPA (\textit{Single-Page Application}) com
  React e Three.js no frontend, e uma API REST com Node.js e MongoDB no
  backend. Contribuições técnicas incluem um mapeador universal de esqueletos
  (\texttt{BoneMapper}), compatível com rigs Mixamo, VRM e CC3; um controlador
  de sincronização labial (\texttt{LipSyncController}) baseado em visemas
  gerados pelo Azure Cognitive Services; e suporte ao padrão VRM Animation
  (\texttt{.vrma}) via \texttt{@pixiv/three-vrm-animation}.

  A avaliação da plataforma foi conduzida por meio do método
  \textit{Expert Walkthrough}, com [N] especialistas nas áreas de IHC,
  Realidade Aumentada e Educação. Os resultados indicam [síntese dos
  resultados — a preencher após coleta]. A plataforma demonstra viabilidade
  técnica e potencial de uso no contexto da Educação~4.0, contribuindo para
  a democratização da autoria de experiências imersivas educacionais.

  \vspace{\onelineskip}
  \noindent\textbf{Palavras-chave:}
  Realidade Aumentada.
  WebAR.
  Narradores Virtuais.
  Avatares 3D.
  Síntese de Voz.
  Lip Sync.
  Plataforma No-Code.
  Educação.
\end{resumo}

% Abstract (inglês)
\begin{resumo}[Abstract]
  \begin{otherlanguage*}{english}
    Augmented Reality (AR) has been widely studied as a pedagogical tool
    capable of promoting interactive and multimodal learning experiences.
    However, the creation of personalized AR content still imposes significant
    technical barriers for educators without a background in software
    development or 3D modeling. A prior exploratory study conducted as part of
    this research identified that existing tools operate in a fragmented
    manner, requiring users to master multiple independent platforms to
    complete a full AR authoring workflow.

    To address this gap, this work presents \textbf{ContAR}, a no-code web
    platform for authoring 3D virtual narrators in educational Augmented
    Reality. The platform integrates, within a single interface, the creation
    and customization of 3D avatars, synthetic speech generation with
    precise lip synchronization via visemes, sequential scene and story
    composition, and browser-accessible AR publication (WebXR Surface AR
    and AR.js Marker AR).

    The architecture follows a Single-Page Application (SPA) model using
    React and Three.js on the frontend, and a REST API with Node.js and
    MongoDB on the backend. Technical contributions include a universal
    skeleton mapper (\texttt{BoneMapper}) compatible with Mixamo, VRM, and
    CC3 rigs; a lip synchronization controller (\texttt{LipSyncController})
    based on visemes generated by Azure Cognitive Services; and support for
    the VRM Animation standard (\texttt{.vrma}) via
    \texttt{@pixiv/three-vrm-animation}.

    Platform evaluation was conducted using the Expert Walkthrough method
    with [N] specialists in HCI, Augmented Reality, and Education. Results
    indicate [summary of results — to be completed after data collection].
    The platform demonstrates technical viability and potential for use in
    the context of Education 4.0, contributing to the democratization of
    immersive educational experience authoring.

    \vspace{\onelineskip}
    \noindent\textbf{Keywords:}
    Augmented Reality.
    WebAR.
    Virtual Narrators.
    3D Avatars.
    Speech Synthesis.
    Lip Sync.
    No-Code Platform.
    Education.
  \end{otherlanguage*}
\end{resumo}

% Listas
\pdfbookmark[0]{\listfigurename}{lof}
\listoffigures*
\cleardoublepage

\pdfbookmark[0]{\listtablename}{lot}
\listoftables*
\cleardoublepage

% Lista de abreviaturas
\begin{siglas}
  \item[ABNT]  Associação Brasileira de Normas Técnicas
  \item[API]   Application Programming Interface
  \item[AR]    Augmented Reality (Realidade Aumentada)
  \item[AR.js] Augmented Reality JavaScript Library
  \item[CRUD]  Create, Read, Update, Delete
  \item[CSS]   Cascading Style Sheets
  \item[GLB]   GL Binary (formato binário do glTF)
  \item[glTF]  Graphics Language Transmission Format
  \item[HTML]  HyperText Markup Language
  \item[HTTP]  HyperText Transfer Protocol
  \item[HTTPS] HTTP Secure
  \item[HCI]   Human-Computer Interaction
  \item[IHC]   Interação Humano-Computador
  \item[JWT]   JSON Web Token
  \item[LLM]   Large Language Model
  \item[MR]    Mixed Reality (Realidade Mista)
  \item[NLP]   Natural Language Processing
  \item[RA]    Realidade Aumentada
  \item[REST]  Representational State Transfer
  \item[RV]    Realidade Virtual
  \item[SDK]   Software Development Kit
  \item[SPA]   Single-Page Application
  \item[SMTP]  Simple Mail Transfer Protocol
  \item[SUS]   System Usability Scale
  \item[TCC]   Trabalho de Conclusão de Curso
  \item[TTS]   Text-to-Speech (Síntese de Voz)
  \item[UFPE]  Universidade Federal de Pernambuco
  \item[VR]    Virtual Reality (Realidade Virtual)
  \item[VRM]   Virtual Reality Model (formato de avatar 3D humanóide)
  \item[VRMA]  VRM Animation (padrão de animação para avatares VRM)
  \item[WebAR] Web-based Augmented Reality
  \item[WebXR] Web Extended Reality API
  \item[XR]    Extended Reality (Realidade Estendida)
\end{siglas}

% Sumário
\pdfbookmark[0]{\contentsname}{toc}
\tableofcontents*
\cleardoublepage

% =============================================================================
% Elementos textuais
% =============================================================================
\textual

% =============================================================================
\chapter{Introdução}
\label{cap:introducao}
% =============================================================================

A Realidade Aumentada (RA) tem se consolidado como uma das tecnologias mais
promissoras para aplicações educacionais. Ao sobrepor elementos virtuais ao
mundo físico, a RA possibilita experiências de aprendizagem interativas,
contextualizadas e multimodais — características frequentemente associadas a
maiores índices de engajamento e retenção de conteúdo
\cite{fonnet2021review, munoz2021education}. Desde as formulações seminal de
Azuma \cite{azuma1997} e a taxonomia de Realidade Mista proposta por Milgram
e Kishino \cite{milgram1994}, o campo evoluiu de sistemas restritos a
laboratórios especializados para experiências acessíveis diretamente pelo
navegador de um smartphone.

Esse avanço foi impulsionado, em grande parte, pela consolidação das
tecnologias WebAR — especialmente a API WebXR \cite{webxr2020} e bibliotecas
como AR.js \cite{arjs} — que eliminam a necessidade de instalação de
aplicativos nativos e permitem que experiências imersivas sejam publicadas,
compartilhadas e acessadas com um simples link. Para contextos educacionais,
onde disponibilidade de dispositivos, restrições de infraestrutura e
heterogeneidade de hardware são fatores críticos, essa acessibilidade
representa uma mudança de paradigma.

Paralelamente, o surgimento de ferramentas de Inteligência Artificial (IA)
para geração de modelos tridimensionais — como Meshy AI \cite{meshy2024},
Luma AI \cite{lumaai2024} e plataformas de criação de avatares como
Ready Player Me \cite{readyplayerme} e Avaturn — democratizou o acesso à
criação de conteúdo 3D, tornando-a viável mesmo para usuários sem formação
em modelagem. Ao mesmo tempo, avanços em síntese de voz neural (TTS) e em
técnicas de sincronização labial via visemas tornaram possível a criação de
narradores virtuais expressivos e naturais.

\section{Motivação e Contexto}
\label{sec:motivacao}

Um estudo exploratório conduzido como etapa anterior a este trabalho
\cite{freitas2025benchmark} investigou a viabilidade de usuários não técnicos
criarem experiências educacionais em RA utilizando ferramentas de IA
disponíveis no mercado. Foram avaliadas três plataformas representativas —
Ready Player Me, Meshy AI e Luma AI — segundo critérios técnicos como
facilidade de uso, qualidade visual, tamanho de arquivo, contagem de
polígonos e compatibilidade com WebAR.

Os resultados desse estudo revelaram que, embora já seja possível criar
modelos 3D de qualidade aceitável para WebAR, o ecossistema atual opera
de forma \textbf{fragmentada}: o usuário precisa navegar por múltiplas
plataformas independentes para concluir um fluxo completo de autoria,
passando pela geração do avatar, conversão de formato, configuração de
áudio, inserção em cena e publicação em RA. Cada transição entre ferramentas
representa uma barreira técnica adicional, especialmente para professores
e educadores sem formação computacional.

A conclusão do estudo anterior apontou diretamente para a necessidade de
\textit{``uma plataforma integrada voltada a usuários não técnicos, capaz
de unificar a criação de avatares, a geração de marcadores, a inserção
de áudio e interações, e a publicação automática em WebAR''}.
Este Trabalho de Conclusão de Curso é a resposta concreta a essa lacuna.

\section{Problema de Pesquisa}
\label{sec:problema}

Diante do cenário descrito, este trabalho endereça o seguinte problema:

\begin{quote}
  \textit{Como uma plataforma web no-code e integrada pode viabilizar a
  autoria de narradores virtuais 3D em Realidade Aumentada por usuários
  sem conhecimento técnico especializado, reduzindo a fragmentação do
  ecossistema atual e tornando o fluxo completo — da criação do avatar
  à publicação em RA — acessível no navegador?}
\end{quote}

\section{Objetivos}
\label{sec:objetivos}

\subsection{Objetivo Geral}
\label{subsec:obj-geral}

Desenvolver, documentar e avaliar o \textbf{ContAR}, uma plataforma web
no-code para autoria de narradores virtuais 3D em Realidade Aumentada
educacional, integrando em um único fluxo a criação de avatares, síntese
de voz com lip sync preciso, composição de cenas e publicação em WebAR
acessível por navegador.

\subsection{Objetivos Específicos}
\label{subsec:obj-especificos}

\begin{enumerate}
  \item Projetar e implementar uma arquitetura full-stack (React + Three.js
        + Node.js + MongoDB) capaz de suportar o ciclo completo de autoria
        de experiências em RA;
  \item Desenvolver um \textit{BoneMapper} universal que permita a
        compatibilidade de animações e poses procedurais com diferentes
        padrões de rig (Mixamo, VRM, CC3 e genérico);
  \item Implementar um sistema de sincronização labial baseado em visemas
        gerados pela API Azure Cognitive Services Speech SDK e pelo padrão
        VRM Animation (\texttt{.vrma});
  \item Integrar modalidades complementares de Realidade Aumentada:
        WebXR Surface AR (posicionamento em superfície) e AR.js Marker AR
        (marcadores fiduciais impressos);
  \item Avaliar a usabilidade da plataforma por meio do método
        \textit{Expert Walkthrough}, identificando barreiras cognitivas,
        pontos de confusão e oportunidades de melhoria no fluxo de autoria;
  \item Analisar a viabilidade do ContAR como ferramenta de apoio à
        Educação~4.0, com foco em professores de educação básica sem
        formação técnica em TI.
\end{enumerate}

\section{Justificativa}
\label{sec:justificativa}

A relevância deste trabalho se articula em três dimensões complementares.

\textbf{Dimensão pedagógica.} Pesquisas em Tecnologia Educacional demonstram
que a RA potencializa o aprendizado ao tornar conceitos abstratos tangíveis,
promovendo construção ativa de conhecimento
\cite{fonnet2021review, munoz2021education}. No Brasil, a adoção de RA
em sala de aula permanece restrita a iniciativas pontuais, em grande parte
pela ausência de ferramentas acessíveis a professores sem formação técnica.
O ContAR endereça diretamente essa barreira ao eliminar a necessidade de
conhecimento em programação ou modelagem 3D.

\textbf{Dimensão tecnológica.} A integração de múltiplos subsistemas
complexos — renderização 3D em tempo real via Three.js, síntese de voz
neural com sincronização labial por visemas, suporte a múltiplos padrões
de avatar (GLB, VRM) e publicação em RA no navegador — em uma única
plataforma coesa representa uma contribuição técnica não trivial, com
potencial de referência para projetos similares na área de WebXR.

\textbf{Dimensão social.} A democratização da autoria de conteúdo
imersivo tem implicações diretas para a equidade educacional. Ao tornar
acessível a criação de narradores virtuais personalizados, o ContAR
abre espaço para a representação de diversidade cultural e identitária
no material didático digital — incluindo aplicações em educação antirracista
e valorização de culturas locais —, perspectiva alinhada a diretrizes
pedagógicas contemporâneas \cite{lgbtqieducation}.

\section{Estrutura do Trabalho}
\label{sec:estrutura}

Este trabalho está organizado em sete capítulos:

O \autoref{cap:referencial} apresenta o referencial teórico, abordando
os fundamentos de Realidade Aumentada, WebAR, avatares 3D nos formatos
GLB e VRM, síntese de voz e lip sync, e plataformas no-code para autoria
de conteúdo educacional.

O \autoref{cap:pesquisa-anterior} descreve a pesquisa exploratória anterior,
cujos resultados motivaram o desenvolvimento do ContAR, incluindo o benchmark
técnico das plataformas de geração 3D avaliadas.

O \autoref{cap:metodologia} apresenta a metodologia de pesquisa adotada,
classificando o trabalho como Pesquisa-Ação combinada com Estudo de Caso,
e descreve o protocolo do Expert Walkthrough utilizado na avaliação.

O \autoref{cap:desenvolvimento} documenta o desenvolvimento da plataforma
ContAR, detalhando as decisões de arquitetura e engenharia, os módulos
técnicos implementados e o processo de design da experiência do usuário.

O \autoref{cap:avaliacao} apresenta os resultados da avaliação do Expert
Walkthrough, incluindo análise das heurísticas de usabilidade, pontuações
SUS e discussão dos achados qualitativos.

O \autoref{cap:conclusao} sintetiza as contribuições do trabalho, discute
as limitações identificadas e aponta direções para trabalhos futuros,
incluindo a proposta de evolução para um sistema de narradores contextuais
com IA integrada.


% =============================================================================
\chapter{Referencial Teórico}
\label{cap:referencial}
% =============================================================================

Este capítulo apresenta os fundamentos teóricos que sustentam o desenvolvimento
e a avaliação do ContAR. As seções estão organizadas de forma a construir
progressivamente o arcabouço conceitual necessário: partindo dos fundamentos
de Realidade Aumentada e WebAR, passando pelos padrões de representação de
avatares 3D, pelos sistemas de síntese de voz e sincronização labial, e
chegando ao conceito de plataformas no-code para autoria de conteúdo educacional.

% ---------------------------------------------------------------------------
\section{Realidade Aumentada: Fundamentos e Evolução}
\label{sec:ra-fundamentos}
% ---------------------------------------------------------------------------

A Realidade Aumentada foi formalmente definida por Azuma \cite{azuma1997}
como um sistema que (1) combina elementos reais e virtuais, (2) é interativo
em tempo real e (3) está registrado em três dimensões. Essa definição,
embora formulada em 1997, permanece como referência na literatura por sua
precisão e abrangência.

A taxonomia proposta por Milgram e Kishino \cite{milgram1994} — o
\textit{Continuum de Realidade Mista} — posiciona a RA entre o ambiente
puramente real e o ambiente puramente virtual. Nesse espectro, a RA
se situa próxima à extremidade real: objetos virtuais são sobrepostos
ao mundo físico, mantendo o usuário ancorado na realidade. Essa propriedade
distingue a RA da Realidade Virtual (RV), onde o usuário é completamente
imerso em um ambiente sintético.

Do ponto de vista técnico, sistemas de RA requerem três componentes
fundamentais: (i) captura do ambiente real (geralmente via câmera),
(ii) rastreamento da posição e orientação do usuário, e
(iii) composição e renderização dos elementos virtuais sobre o feed real.
As estratégias de rastreamento evoluíram de marcadores fiduciais
(padrões impressos de alto contraste) para rastreamento markerless baseado
em visão computacional (SLAM — \textit{Simultaneous Localization and Mapping}),
e mais recentemente para sistemas híbridos que combinam ambas as abordagens.

\subsection{Realidade Aumentada na Educação}
\label{subsec:ra-educacao}

O potencial pedagógico da RA tem sido amplamente investigado. Revisões
sistemáticas da literatura, como as conduzidas por Fonnet e Prie
\cite{fonnet2021review} e Muñoz-Cristóbal et al. \cite{munoz2021education},
demonstram que a RA pode promover ganhos significativos de engajamento,
motivação intrínseca e compreensão conceitual, especialmente em disciplinas
que se beneficiam de visualização espacial — como Ciências, Geografia,
Biologia e Matemática.

Contudo, esses estudos também evidenciam que a maioria das aplicações
educacionais de RA é desenvolvida por equipes especializadas, com baixa
participação ativa de professores no processo de criação. Isso contrasta
com princípios contemporâneos de design instrucional, que enfatizam a
importância da autoria docente para a contextualização pedagógica do
material \cite{ruiz2020arbeginner}. O ContAR endereça diretamente essa
limitação ao colocar a capacidade de autoria nas mãos do educador.

% ---------------------------------------------------------------------------
\section{WebAR: Realidade Aumentada no Navegador}
\label{sec:webar}
% ---------------------------------------------------------------------------

O paradigma WebAR emergiu como alternativa às aplicações nativas de RA,
eliminando a necessidade de instalação de software e aproveitando a
ubiquidade dos navegadores modernos. A transição foi viabilizada pelo
avanço da API WebXR \cite{webxr2020}, que padronizou o acesso a recursos
de XR nos navegadores, e pelo amadurecimento de bibliotecas JavaScript
especializadas.

\subsection{API WebXR}
\label{subsec:webxr}

A WebXR Device API é uma especificação do W3C que fornece acesso
padronizado a dispositivos de Realidade Virtual e Aumentada diretamente
no navegador. Ela substitui as APIs legadas WebVR e A-Frame AR, oferecendo
suporte a recursos avançados como rastreamento de superfície (\textit{hit-test}),
oclusão de profundidade e estimativa de iluminação ambiente.

No contexto do ContAR, a WebXR é utilizada para o modo Surface AR:
o usuário aponta a câmera para uma superfície plana (piso, mesa), o
sistema detecta automaticamente a superfície via rastreamento SLAM do
dispositivo, e o narrador virtual é posicionado sobre ela com ancoragem
estável. Esse modo dispensa qualquer marcador impresso e funciona em
dispositivos Android com Chrome e iOS com Safari (suporte parcial).

\subsection{AR.js e Marcadores Fiduciais}
\label{subsec:arjs}

O AR.js \cite{arjs} é uma biblioteca JavaScript de código aberto para RA
baseada em marcadores fiduciais. Originalmente construída sobre A-Frame e
ARToolKit, o AR.js utiliza visão computacional para detectar padrões
impressos (marcadores) no feed da câmera e posicionar modelos 3D sobre
eles com rastreamento estável.

As vantagens do AR.js para contextos educacionais são significativas:
baixo custo computacional (funciona em smartphones com 2+ anos),
compatibilidade ampla com navegadores (Chrome, Firefox, Safari),
e a possibilidade de criar materiais didáticos impressos —
cartilhas, pôsteres, fichas — que servem como âncoras para os
narradores virtuais. No ContAR, o modo Marker AR usa AR.js via um iframe
leve, permitindo que histórias criadas na plataforma sejam visualizadas
sobre marcadores Hiro (padrão) ou customizados.

\subsection{Three.js como Motor de Renderização}
\label{subsec:threejs}

O Three.js \cite{threejs} é a biblioteca JavaScript de renderização 3D
mais utilizada na web, com mais de 97.000 estrelas no GitHub. Baseada em
WebGL, ela abstrai as complexidades da API gráfica de baixo nível e
fornece um sistema de cena completo com câmeras, luzes, materiais,
geometrias, carregadores de modelos e um loop de renderização configurável.

No ContAR, o Three.js r183 é o motor principal da experiência 3D,
responsável pela renderização dos avatares no editor e nas cenas de RA.
A escolha se justifica pela maturidade da biblioteca, pelo suporte nativo
ao formato glTF/GLB via \texttt{GLTFLoader}, pela integração com
\texttt{DRACOLoader} para modelos comprimidos, e pela compatibilidade
com o plugin \texttt{VRMLoaderPlugin} para avatares no padrão VRM.

% ---------------------------------------------------------------------------
\section{Avatares 3D: Formatos GLB e VRM}
\label{sec:avatares}
% ---------------------------------------------------------------------------

A representação de personagens virtuais humanóides em aplicações WebAR
envolve dois aspectos críticos: o formato de armazenamento do modelo 3D
e a estrutura de esqueleto (rig) que define como o personagem pode ser
animado.

\subsection{O Formato glTF e GLB}
\label{subsec:gltf-glb}

O glTF (\textit{Graphics Language Transmission Format}) é um padrão aberto
do Khronos Group para transmissão e carregamento de cenas e modelos 3D.
Projetado para ser o ``JPEG dos modelos 3D'', o glTF prioriza compactação,
velocidade de carregamento e compatibilidade com diferentes engines.
O GLB é a variante binária do glTF, que empacota geometria, texturas e
animações em um único arquivo, ideal para transmissão na web.

Plataformas de criação de avatares como Avaturn, Ready Player Me e
CharacterStudio utilizam o GLB como formato de exportação padrão.
O arquivo inclui a malha poligonal do personagem, materiais PBR
(\textit{Physically Based Rendering}), esqueleto de ossos e, em muitos
casos, clips de animação embutidos (idle, walk, etc.).

\subsection{O Padrão VRM}
\label{subsec:vrm}

O VRM (\textit{Virtual Reality Model}) é uma extensão do glTF desenvolvida
pela Pixiv Inc. especificamente para avatares humanóides. Além da geometria
e animação padrão do glTF, o VRM define metadados estruturados que incluem:
(i) mapeamento humanóide de ossos (\textit{humanoid bone names}),
padronizando denominações como \texttt{hips}, \texttt{head},
\texttt{leftUpperArm}; (ii) \textit{blend shapes} para expressões faciais
(feliz, triste, surpreso, piscar) e visemas de fala; (iii) metadados de
autoria e licença; e (iv) física de cabelo e roupas via \textit{spring bones}.

A padronização do VRM tem implicações importantes para a animação e o
lip sync: ao definir nomes canônicos para ossos e blend shapes, permite
que animações, expressões e visemas criados para qualquer avatar VRM
funcionem corretamente em qualquer outro avatar VRM compatível. O ContAR
explora essa propriedade para aplicar poses procedurais, visemas de fala
e animações no formato VRM Animation (\texttt{.vrma}) de forma
independente do modelo específico.

\subsection{BoneMapper: Mapeamento Universal de Esqueletos}
\label{subsec:bonemapper}

Um desafio central no desenvolvimento do ContAR foi garantir a
compatibilidade de animações e poses com avatares de diferentes origens
(Avaturn, VRoid, CharacterStudio, modelos genéricos). Cada ecossistema
usa convenções distintas de nomenclatura de ossos: Avaturn/Mixamo usa
\texttt{mixamorigHips}; VRM usa \texttt{J\_Bip\_C\_Hips}; CC3 usa
\texttt{CC\_Base\_Hip}; modelos genéricos podem usar quaisquer nomes.

O \texttt{BoneMapper} desenvolvido para o ContAR resolve esse problema
identificando automaticamente o padrão de rig do avatar carregado
(via análise heurística dos nomes de ossos) e mantendo um mapeamento
interno de nomes canônicos para referências de objetos Three.js. Isso
permite que todo o sistema de poses procedurais e animações opere
sobre nomes abstratos (\texttt{leftUpperArm}, \texttt{hips}, etc.),
independentemente da nomenclatura específica do modelo.

% ---------------------------------------------------------------------------
\section{Síntese de Voz e Sincronização Labial}
\label{sec:tts-lipsync}
% ---------------------------------------------------------------------------

Narradores virtuais expressivos requerem não apenas voz sintética de
qualidade, mas também sincronização precisa dos movimentos labiais com
o áudio gerado — o que a literatura denomina \textit{lip sync}.

\subsection{Text-to-Speech Neural}
\label{subsec:tts}

A geração de fala a partir de texto (\textit{Text-to-Speech}, TTS) evoluiu
radicalmente com modelos neurais profundos. Sistemas como os disponíveis
no Azure Cognitive Services Speech SDK \cite{azurespeech} utilizam redes
neurais convolucionais e transformers para gerar voz que, em avaliações
de escores MOS (\textit{Mean Opinion Score}), se aproxima da qualidade da
fala humana natural. O Azure Neural TTS oferece mais de 400 vozes em
+140 idiomas e localidades, incluindo variantes brasileiras do português
(\texttt{pt-BR-FranciscaNeural}, \texttt{pt-BR-AntonioNeural}, entre outras).

\subsection{Visemas e Sincronização Labial}
\label{subsec:visemas}

Um \textit{visema} é a representação visual de um fonema — o conjunto de
posições labiais e faciais que corresponde a um som específico da fala.
O Azure Speech SDK fornece eventos \texttt{visemeReceived} durante a
síntese de voz, entregando identificadores de visema (IDs 0--21, cobrindo
os principais fonemas do inglês e do português) com timestamps de precisão
de 100 nanosegundos.

No ContAR, esses eventos são capturados no backend e serializados como
uma \textit{timeline de visemas}: uma lista ordenada de objetos
\texttt{\{start, end, value\}}, onde \texttt{value} é um código Rhubarb
(A, B, C, D, E, F, G, H, X) mapeado a partir do ID de visema Azure.
Essa timeline é transmitida ao frontend junto com o áudio base64 e
consumida pelo \texttt{LipSyncController} para animar, frame a frame,
os blend shapes de boca do avatar com crossfade suavizado.

\subsection{LipSyncController}
\label{subsec:lipsync-controller}

O \texttt{LipSyncController} implementado no ContAR gerencia a animação
dos blend shapes de fala de forma independente do rig do avatar, operando
sobre grupos semânticos de alvos: \texttt{aa} (boca aberta), \texttt{oh}
(boca arredondada), \texttt{ee} (boca esticada), \texttt{fv} (dentes sobre
lábio inferior) e \texttt{mbp} (lábios fechados). Para cada grupo, o
controlador identifica automaticamente os blend shapes correspondentes no
modelo via padrões de nomenclatura (p.ex., \texttt{viseme\_aa},
\texttt{jawOpen}, \texttt{mouthSmile}).

O sistema opera em dois modos: (1) modo \textit{timeline}, que usa a
sequência de visemas Azure com interpolação temporal; e
(2) modo \textit{heurístico}, que estima o movimento labial em tempo real
a partir da análise de amplitude e frequência do sinal de áudio via Web
Audio API — útil quando não há timeline de visemas disponível
(Web Speech API, arquivos de áudio importados pelo usuário).

% ---------------------------------------------------------------------------
\section{Plataformas No-Code para Autoria de Conteúdo Educacional}
\label{sec:nocode}
% ---------------------------------------------------------------------------

O conceito de \textit{no-code} refere-se a ferramentas que permitem a
criação de sistemas ou conteúdo funcional sem a necessidade de escrever
código de programação, geralmente por meio de interfaces visuais,
configurações dirigidas por formulários ou composição baseada em componentes
\cite{nocode2022}. No contexto educacional, plataformas no-code são
especialmente relevantes por expandir a capacidade de criação a professores
e educadores que possuem domínio pedagógico mas não técnico.

Ferramentas como Google Sites, Padlet e Canva demonstraram o potencial
do modelo no-code para empoderar criadores de conteúdo não técnicos.
No domínio específico da RA educacional, entretanto, as opções no-code
permanecem limitadas: a maioria das ferramentas exige pelo menos noções
de desenvolvimento web, uso de motores gráficos (Unity, Unreal Engine)
ou dependência de serviços proprietários com curvas de aprendizado
significativas.

O ContAR posiciona-se nessa lacuna ao oferecer um fluxo completo de autoria
— da criação do avatar à publicação em RA — inteiramente dentro do navegador,
sem que o usuário precise instalar software, escrever código ou gerenciar
arquivos entre múltiplas plataformas. A comparação com ferramentas
equivalentes é aprofundada no \autoref{cap:pesquisa-anterior}.


% =============================================================================
\chapter{Pesquisa Anterior: Benchmark de Ferramentas de Criação 3D para WebAR}
\label{cap:pesquisa-anterior}
% =============================================================================

Este capítulo apresenta a pesquisa exploratória conduzida como etapa
anterior ao desenvolvimento do ContAR. Publicada como artigo científico
\cite{freitas2025benchmark}, essa pesquisa identificou a lacuna que
motivou o presente trabalho.

\section{Contexto e Motivação}
\label{sec:benchmark-contexto}

O estudo foi conduzido a partir da seguinte questão: até que ponto usuários
sem conhecimento técnico conseguem criar experiências educacionais completas
em RA utilizando ferramentas modernas de IA e tecnologias WebAR? Para
respondê-la, foi adotada uma metodologia de mapeamento exploratório combinado
com benchmark técnico comparativo.

\section{Ferramentas Avaliadas}
\label{sec:benchmark-ferramentas}

Três plataformas foram selecionadas por representarem diferentes paradigmas
de criação 3D:

\begin{itemize}
  \item \textbf{Ready Player Me}: criação de avatares humanos estilizados
        com foco em acessibilidade, via interface guiada por selfie ou
        configuração manual de atributos físicos.
  \item \textbf{Meshy AI}: geração de modelos 3D a partir de descrições
        textuais (text-to-3D) ou imagens de referência, com ênfase em
        otimização automática de malhas.
  \item \textbf{Luma AI (Genie)}: geração neural de objetos tridimensionais
        com foco em fidelidade visual e realismo fotográfico.
\end{itemize}

\section{Critérios e Resultados}
\label{sec:benchmark-resultados}

Os modelos gerados foram avaliados segundo sete critérios técnicos:
facilidade de uso, qualidade visual, processo de exportação GLB, tamanho
do arquivo, contagem de polígonos, compatibilidade com WebAR e tempo total
do fluxo. Cada critério foi pontuado em escala ordinal de 0 a 3 (máximo: 21
pontos). Os resultados são sintetizados na Tabela~\ref{tab:benchmark-anterior}.

\begin{table}[ht]
  \caption{Síntese do benchmark técnico das plataformas de criação 3D.}
  \label{tab:benchmark-anterior}
  \centering
  \begin{tabular}{lrrrrrrrc}
    \toprule
    \textbf{Plataforma} & \textbf{A} & \textbf{B} & \textbf{C} &
    \textbf{D} & \textbf{E} & \textbf{F} & \textbf{G} & \textbf{Total} \\
    \midrule
    Ready Player Me & 3 & 2 & 3 & 2 & 2 & 3 & 3 & \textbf{18} \\
    Meshy AI        & 3 & 3 & 3 & 3 & 3 & 3 & 2 & \textbf{20} \\
    Luma AI         & 2 & 3 & 2 & 1 & 1 & 2 & 2 & \textbf{13} \\
    \bottomrule
  \end{tabular}
  \vspace{0.3em}
  \footnotesize{A=Facilidade de uso; B=Qualidade visual; C=Exportação GLB;
  D=Tamanho do arquivo; E=Polígonos; F=Compatibilidade WebAR; G=Tempo do fluxo.}
\end{table}

O Meshy AI obteve a maior pontuação geral (20/21), destacando-se pela
otimização técnica dos modelos. O Ready Player Me obteve segunda colocação
(18/21), com vantagem significativa na facilidade de uso e no tempo total
do fluxo. A Luma AI, embora com excelente qualidade visual, apresentou
modelos com tamanho de arquivo e contagem de polígonos excessivos para
dispositivos móveis (13/21).

\section{A Lacuna Identificada}
\label{sec:lacuna}

A conclusão central do estudo anterior foi que, apesar da qualidade
individual das ferramentas, o ecossistema opera de forma fragmentada.
O usuário precisa: (1) criar o avatar em uma plataforma; (2) exportar
em formato GLB; (3) validar o arquivo externamente; (4) inserir em uma
cena WebAR; (5) configurar áudio separadamente; (6) publicar em um servidor.
Cada etapa exige decisões técnicas que representam barreiras para
usuários não técnicos.

Essa fragmentação é a motivação direta para o desenvolvimento do ContAR,
que integra todas essas etapas em um único fluxo unificado dentro do
navegador.


% =============================================================================
\chapter{Metodologia da Pesquisa}
\label{cap:metodologia}
% =============================================================================

\section{Classificação da Pesquisa}
\label{sec:classificacao}

Este trabalho é classificado como \textbf{Pesquisa-Ação} combinada com
\textbf{Estudo de Caso}. A Pesquisa-Ação caracteriza-se pela participação
ativa do pesquisador no fenômeno estudado, com iterações cíclicas entre
observação, ação e reflexão \cite{tripp2005}. Essa classificação reflete
a natureza do trabalho: o pesquisador não apenas observou o problema da
fragmentação de ferramentas de RA, mas construiu e avaliou ativamente uma
solução para ele.

O Estudo de Caso é adotado na fase de avaliação, onde o ContAR é investigado
como objeto de análise em profundidade, com foco nos processos de uso e
nas percepções dos especialistas participantes.

\section{Etapas da Pesquisa}
\label{sec:etapas}

A pesquisa foi estruturada em quatro etapas sequenciais:

\begin{enumerate}
  \item \textbf{Mapeamento exploratório e benchmark}: levantamento e avaliação
        técnica de ferramentas de geração 3D para WebAR (detalhado no
        \autoref{cap:pesquisa-anterior}), que identificou a lacuna que
        motivou o desenvolvimento do ContAR.
  \item \textbf{Design e desenvolvimento}: projeto da arquitetura e
        implementação iterativa da plataforma ContAR, documentada no
        \autoref{cap:desenvolvimento}.
  \item \textbf{Avaliação}: condução do Expert Walkthrough com especialistas,
        coleta de dados e análise, apresentada no \autoref{cap:avaliacao}.
  \item \textbf{Síntese e publicação}: consolidação dos resultados em forma
        de relatório acadêmico (este TCC) e artigo científico submetido ao
        SVR \cite{freitas2025contar}.
\end{enumerate}

\section{Expert Walkthrough: Protocolo de Avaliação}
\label{sec:expert-walkthrough}

O método \textit{Expert Walkthrough} (também conhecido como
\textit{Cognitive Walkthrough} especializado) é uma técnica de avaliação
de usabilidade na qual especialistas percorrem sistematicamente as tarefas
de um sistema, verificando cada passo segundo heurísticas de usabilidade
e registrando problemas, dificuldades e sugestões de melhoria
\cite{lewis1990walkthrough, nielsen1994heuristics}.

A escolha desse método se justifica por: (i) ser adequado para a fase
de avaliação formativa de protótipos; (ii) exigir de 3 a 5 avaliadores
(escala viável para um TCC individual); (iii) identificar cerca de 75\%
dos problemas de usabilidade em estudos empíricos \cite{nielsen1994heuristics};
e (iv) permitir a coleta de dados ricos e contextualizados sobre a
experiência de especialistas com o sistema.

\subsection{Participantes}
\label{subsec:participantes}

Foram recrutados [N=?] especialistas com perfis complementares, seguindo
a recomendação de Nielsen \cite{nielsen1994heuristics} de 3 a 5 avaliadores
com conhecimento em usabilidade:

\begin{itemize}
  \item [Especialista 1] — Perfil: [IHC/Educação/XR]
  \item [Especialista 2] — Perfil: [IHC/Educação/XR]
  \item [Especialista 3] — Perfil: [IHC/Educação/XR]
  % Adicionar participantes após coleta
\end{itemize}

\subsection{Tarefas}
\label{subsec:tarefas}

Foram definidas cinco tarefas cobrindo o fluxo completo de autoria:

\begin{description}
  \item[T1 — Criar e carregar um avatar:]
    Acessar a plataforma, criar ou selecionar um avatar usando uma das
    opções disponíveis (Avaturn, galeria CC0, upload GLB/VRM) e
    visualizá-lo no editor 3D.
  \item[T2 — Configurar a fala e gerar voz:]
    Digitar um texto no campo de fala do narrador, selecionar uma voz
    e gerar o áudio com sincronização labial via TTS.
  \item[T3 — Definir pose e salvar a cena:]
    Escolher uma pose no seletor, definir um título para a cena e
    salvá-la no servidor.
  \item[T4 — Compor e publicar uma história:]
    Adicionar múltiplas cenas a uma história, atribuir título e
    descrição, e obter o link de compartilhamento público.
  \item[T5 — Visualizar em Realidade Aumentada:]
    Acessar o link de AR no celular, posicionar o narrador em uma
    superfície real (WebXR) ou sobre um marcador impresso (AR.js),
    e reproduzir a história com áudio.
\end{description}

\subsection{Heurísticas Aplicadas}
\label{subsec:heuristicas}

A avaliação foi baseada nas 10 heurísticas de usabilidade de Nielsen
\cite{nielsen1994heuristics}, complementadas por heurísticas específicas
para sistemas de Realidade Estendida propostas por Gabbard et al.
\cite{gabbard2006xr}:

\begin{enumerate}
  \item Visibilidade do status do sistema
  \item Correspondência entre o sistema e o mundo real
  \item Controle e liberdade do usuário
  \item Consistência e padrões
  \item Prevenção de erros
  \item Reconhecimento em vez de lembrança
  \item Flexibilidade e eficiência de uso
  \item Estética e design minimalista
  \item Ajuda ao usuário para reconhecer, diagnosticar e recuperar erros
  \item Ajuda e documentação
  \item[XR1.] Estabilidade e ancoragem de objetos virtuais
  \item[XR2.] Feedback de rastreamento e estado do AR
  \item[XR3.] Ergonomia da interação em ambiente móvel
\end{enumerate}

\subsection{Instrumentos de Coleta}
\label{subsec:instrumentos}

Os dados foram coletados por meio de três instrumentos:

\begin{itemize}
  \item \textbf{Protocolo Think Aloud}: os especialistas verbalizaram
        seus pensamentos durante a execução das tarefas, registrados
        em áudio com consentimento.
  \item \textbf{Ficha de registro de problemas}: formulário estruturado
        para catalogar problemas por heurística violada, severidade
        (escala de Nielsen 0--4) e localização na interface.
  \item \textbf{SUS (System Usability Scale)} \cite{brooke1996sus}:
        questionário de 10 itens em escala Likert aplicado após a
        sessão, gerando um score de 0 a 100.
\end{itemize}


% =============================================================================
\chapter{ContAR: Desenvolvimento da Plataforma}
\label{cap:desenvolvimento}
% =============================================================================

\section{Visão Geral da Arquitetura}
\label{sec:arquitetura}

O ContAR é uma aplicação web full-stack organizada em duas camadas
principais: um frontend SPA (\textit{Single-Page Application}) e uma
API REST no backend. Ambas são implantadas como serviços independentes
em contêineres Docker, comunicando-se via HTTPS com suporte a proxy
reverso nginx.

\begin{figure}[ht]
  \centering
  % \includegraphics[width=0.9\textwidth]{figuras/arquitetura-geral.png}
  \fbox{\parbox{0.85\textwidth}{\centering\vspace{2cm}
    [Diagrama de arquitetura — a inserir]
  \vspace{2cm}}}
  \caption{Arquitetura geral da plataforma ContAR.}
  \label{fig:arquitetura}
\end{figure}

\subsection{Stack Tecnológico}
\label{subsec:stack}

\begin{description}
  \item[Frontend] React 19 + Vite + Tailwind CSS. Three.js r183 para
    renderização 3D. Zustand para gerenciamento de estado com persistência
    via \textit{localStorage}. React Router para navegação SPA.
    i18next para internacionalização (PT, EN, ES, FR).

  \item[Backend] Node.js 22 (LTS) com Express.js. MongoDB com Mongoose
    para persistência de cenas, histórias e usuários. JWT para
    autenticação. Nodemailer (Gmail SMTP) para e-mails transacionais.
    Azure Cognitive Services Speech SDK para síntese de voz e visemas.
    Multer para upload de arquivos GLB/VRM e áudio.

  \item[Infraestrutura] Docker (frontend: \texttt{node:22-alpine} +
    nginx; backend: \texttt{node:22-slim} para compatibilidade com
    dependências glibc do Azure SDK). Deploy na plataforma Render.com.
\end{description}

\section{Módulos do Frontend}
\label{sec:frontend}

\subsection{SceneCanvas: Motor de Renderização 3D}
\label{subsec:scenecanvas}

O \texttt{SceneCanvas} é o componente React central do editor, responsável
por encapsular o ciclo de vida completo da cena Three.js. Suas
responsabilidades incluem:

\begin{itemize}
  \item Carregamento de modelos GLB e VRM via \texttt{GLTFLoader} com
        suporte a \texttt{DRACOLoader} (compressão Draco) e
        \texttt{VRMLoaderPlugin};
  \item Gerenciamento de câmera com \texttt{OrbitControls} e iluminação
        HDR via \texttt{RGBELoader};
  \item Execução do loop de renderização com pausa automática quando o
        componente fica fora da viewport (\texttt{IntersectionObserver});
  \item Integração com o \texttt{AnimationController} para poses
        procedurais e animações;
  \item Integração com o \texttt{LipSyncController} para sincronização
        labial em tempo real;
  \item Painel de ferramentas de desenvolvimento (ativado via
        \texttt{?dev} na URL) com debug de lip sync, catálogo de ossos
        e override manual de mapeamento.
\end{itemize}

\subsection{AnimationController: Poses e Animações}
\label{subsec:animation-controller}

O \texttt{AnimationController} gerencia o \texttt{AnimationMixer} do
Three.js com crossfade suave entre clipes, e implementa micro-animações
procedurais: piscar de olhos com timing aleatório, respiração via
oscilação do osso do tórax, e gestos de locutor (movimentos de braço
sincronizados com o modo \textit{speaker}).

O sistema de poses suporta 16 presets: 5 animadas (idle, walk, run,
dance, speaker) e 11 estáticas (neutral, wave, hands on hips, salute,
arms crossed, think, point, bow, pray, shrug, t-pose). As poses estáticas
são implementadas como rotações procedurais de ossos no sistema de
coordenadas local, usando uma abstração que compensa a orientação do
rig exportado (\texttt{rotateBoneDeg}).

O \texttt{AnimationController} também retargetiza clipes de animação
via o \texttt{BoneMapper}: ao receber um clipe com nomes de ossos
Mixamo ou VRM, renomeia as trilhas para corresponder à nomenclatura do
avatar atualmente carregado, garantindo compatibilidade cruzada.

Animações no padrão VRM Animation (\texttt{.vrma}) são suportadas via
o pacote \texttt{@pixiv/three-vrm-animation}: o arquivo é carregado por
um \texttt{GLTFLoader} com \texttt{VRMAnimationLoaderPlugin}, e a função
\texttt{createVRMAnimationClip} gera um \texttt{AnimationClip} Three.js
mapeado para os ossos humanóides do VRM específico.

\subsection{LipSyncController}
\label{subsec:lipsync-impl}

O \texttt{LipSyncController} abstrai a animação de blend shapes de fala
do avatar, independentemente do rig ou da convenção de nomenclatura.
Ao ser instanciado com um modelo, ele percorre automaticamente todas as
malhas com \textit{morph targets} e os agrupa semanticamente por
padrões de expressão regular (p.ex., \texttt{/viseme.*aa/i},
\texttt{/jawopen/i}, \texttt{/mouthsmile/i}).

Durante o loop de renderização, recebe valores de ativação por grupo
(0.0 a 1.0) e os distribui para os blend shapes correspondentes com
ponderação proporcional. A transição entre frames usa interpolação
linear configurável (\textit{crossfade}), produzindo movimentos labiais
fluidos sem artefatos de pop.

\section{Sistema de Cenas e Histórias}
\label{sec:cenas-historias}

O modelo de dados do ContAR é estruturado em torno de dois conceitos:
\textbf{cenas} e \textbf{histórias}. Uma cena representa um ``quadro''
narrativo completo — avatar, pose, fala, áudio e configuração de transformação.
Uma história é uma sequência ordenada de cenas identificadas por ID,
que o visualizador percorre automaticamente, reproduzindo o áudio de
cada cena e avançando ao término.

Esse modelo permite que um professor componha, de forma não-linear e
iterativa, narrativas complexas: pode criar cenas individualmente,
revisitá-las, reordenar e reagrupar em diferentes histórias, e publicar
com um link único. O visualizador (\texttt{/story/:id}) é público e
não requer autenticação.

\section{Realidade Aumentada no ContAR}
\label{sec:ar-contar}

O ContAR implementa duas modalidades complementares de RA, acessíveis
na rota \texttt{/ar}:

\textbf{Surface AR (WebXR):} usa a API WebXR com suporte a
\textit{hit-test} para detectar superfícies planas via SLAM do
dispositivo. O avatar é posicionado sobre um retículo e ancorado ao
toque do usuário. Suporta controle de escala persistido no
\textit{localStorage} e reprodução automática de histórias com áudio.

\textbf{Marker AR (AR.js):} carrega um iframe com uma cena A-Frame + AR.js
configurada para detectar o marcador Hiro padrão ou marcadores
customizados. O avatar é renderizado sobre o marcador com orientação e
escala configuráveis via parâmetros de URL (\texttt{?modelUrl=},
\texttt{?scale=}).


% =============================================================================
\chapter{Avaliação: Expert Walkthrough}
\label{cap:avaliacao}
% =============================================================================

\section{Condução das Sessões}
\label{sec:conducao}

% [A preencher após coleta de dados]
As sessões foram conduzidas [individualmente / remotamente via videoconferência]
com duração média de [X] minutos cada. Os especialistas foram orientados a
verbalizar seus pensamentos em voz alta (protocolo Think Aloud) durante
a execução das tarefas, e receberam acesso à versão de produção da
plataforma sem assistência do pesquisador, exceto em casos de bloqueio
total na tarefa.

\section{Resultados Quantitativos}
\label{sec:resultados-quant}

\subsection{SUS Score}
\label{subsec:sus}

% [A preencher após coleta de dados]
A pontuação média na System Usability Scale foi de \textbf{[X]} pontos
(desvio padrão: [Y]), o que corresponde à classificação ``[Excellent /
Good / OK]'' segundo a escala de adjunto de Bangor et al. \cite{bangor2008sus}.
A Tabela~\ref{tab:sus} apresenta as pontuações individuais.

\begin{table}[ht]
  \caption{Pontuações SUS individuais.}
  \label{tab:sus}
  \centering
  \begin{tabular}{lc}
    \toprule
    \textbf{Especialista} & \textbf{SUS Score} \\
    \midrule
    E1 & --- \\
    E2 & --- \\
    E3 & --- \\
    \midrule
    \textbf{Média} & \textbf{---} \\
    \bottomrule
  \end{tabular}
\end{table}

\subsection{Problemas Identificados por Heurística}
\label{subsec:problemas}

% [A preencher após análise dos dados]
Foram identificados [N] problemas de usabilidade, distribuídos conforme
a Tabela~\ref{tab:heuristicas}.

\begin{table}[ht]
  \caption{Distribuição de problemas por heurística violada.}
  \label{tab:heuristicas}
  \centering
  \begin{tabular}{lcc}
    \toprule
    \textbf{Heurística} & \textbf{Ocorrências} & \textbf{Severidade Média} \\
    \midrule
    H1 — Visibilidade do status & --- & --- \\
    H3 — Controle e liberdade  & --- & --- \\
    H6 — Reconhecimento        & --- & --- \\
    H9 — Recuperação de erros  & --- & --- \\
    XR2 — Feedback AR          & --- & --- \\
    \bottomrule
  \end{tabular}
\end{table}

\section{Resultados Qualitativos}
\label{sec:resultados-qual}

% [A preencher após análise dos dados]
A análise das verbalizações registradas durante o Think Aloud revelou
[principais achados qualitativos]. Os especialistas identificaram como
pontos fortes [X, Y, Z], e como principais oportunidades de melhoria [A, B, C].

\section{Discussão}
\label{sec:discussao}

% [A preencher após coleta e análise]
Os resultados obtidos permitem [análise em relação à questão de pesquisa].
Em comparação com o benchmark anterior \cite{freitas2025benchmark},
o ContAR demonstra [...]


% =============================================================================
\chapter{Conclusão e Trabalhos Futuros}
\label{cap:conclusao}
% =============================================================================

\section{Síntese das Contribuições}
\label{sec:contribuicoes}

Este trabalho apresentou o ContAR, uma plataforma web no-code para autoria
de narradores virtuais 3D em Realidade Aumentada educacional, respondendo
à lacuna identificada em pesquisa anterior: a ausência de um ecossistema
integrado que permita a usuários não técnicos completar o fluxo de autoria
em RA sem fragmentação entre ferramentas.

As contribuições se articulam em três dimensões:

\textbf{Contribuição técnica:} (i) arquitetura full-stack integrada para
autoria WebAR no-code; (ii) \texttt{BoneMapper} com compatibilidade
universal de rigs (Mixamo, VRM, CC3, genérico); (iii)
\texttt{LipSyncController} com modos timeline (visemas Azure) e heurístico
(análise de áudio); (iv) suporte ao padrão VRM Animation para avatares
VRM; (v) publicação em RA dual-mode (WebXR Surface AR + AR.js Marker AR).

\textbf{Contribuição metodológica:} avaliação de usabilidade via Expert
Walkthrough com heurísticas de Nielsen complementadas por heurísticas XR
específicas, produzindo um protocolo replicável para avaliação de sistemas
de autoria WebAR.

\textbf{Contribuição social:} democratização da autoria de experiências
imersivas educacionais, com implicações para a representação de diversidade
cultural no material didático digital e para a Educação~4.0 no contexto
brasileiro.

\section{Limitações}
\label{sec:limitacoes}

O presente trabalho apresenta algumas limitações que devem ser consideradas
na interpretação dos resultados:

\begin{itemize}
  \item O Expert Walkthrough não substitui testes com usuários finais
        (professores e estudantes). Uma avaliação com o público-alvo primário
        poderá revelar problemas de usabilidade distintos dos identificados
        pelos especialistas.
  \item O suporte WebXR para Surface AR depende da implementação do
        navegador no dispositivo. Dispositivos iOS com versões antigas
        do Safari apresentam suporte limitado.
  \item A síntese de voz via Azure TTS requer configuração de credenciais
        de API, o que introduz uma barreira de configuração para implantações
        independentes.
  \item Os modelos de avatar gerados por IA de terceiros (Avaturn,
        CharacterStudio) estão sujeitos a mudanças nas APIs e políticas
        de serviço dessas plataformas.
\end{itemize}

\section{Trabalhos Futuros}
\label{sec:trabalhos-futuros}

Com base nos resultados e nas limitações identificadas, sugere-se como
direções para trabalhos futuros:

\textbf{Avaliação com usuários finais:} conduzir testes de usabilidade
com professores da educação básica, avaliando o impacto do ContAR na
redução do tempo de criação de material didático imersivo em comparação
com fluxos fragmentados existentes.

\textbf{Narradores contextuais com IA:} integrar um LLM (Large Language
Model) para que o narrador responda a perguntas dos estudantes em tempo
real, transformando o personagem virtual de um apresentador estático em
um agente conversacional educacional. Essa proposta, sugerida pelo
orientador Prof. Dr. [Nome], representa a fronteira de pesquisa mais
promissora e a ponte natural para um projeto de mestrado.

\textbf{Expressões faciais VRM:} explorar o sistema de expressões do
padrão VRM (\texttt{vrm.expressionManager}) para adicionar emoções
sincronizadas ao conteúdo narrativo — felicidade, surpresa, seriedade —
enriquecendo a comunicação paraverbal do narrador.

\textbf{Eye tracking (VRM LookAt):} implementar o sistema
\texttt{vrm.lookAt} para que os olhos do narrador acompanhem a câmera
do usuário, aumentando a sensação de presença e engajamento.

\textbf{Biblioteca de poses e animações:} desenvolver ou curatorial um
repositório de animações Mixamo convertidas para GLB e de arquivos VRMA
compatíveis, expandindo as opções expressivas do narrador sem dependência
de serviços externos.

\textbf{Aplicações em Educação Antirracista:} investigar o potencial do
ContAR para a criação de narradores com representações diversas de etnia,
cultura e identidade, contribuindo para práticas pedagógicas de valorização
da diversidade alinhadas a diretrizes educacionais brasileiras contemporâneas.


% =============================================================================
% Elementos pós-textuais
% =============================================================================
\postextual

% Referências
\bibliographystyle{abntex2-alf}
\bibliography{referencias}

% As referências abaixo devem estar no arquivo referencias.bib
% Exemplos de entradas necessárias:
%
% @inproceedings{azuma1997,
%   author    = {Azuma, Ronald T.},
%   title     = {A Survey of Augmented Reality},
%   journal   = {Presence: Teleoperators and Virtual Environments},
%   year      = {1997}, volume = {6}, number = {4}, pages = {355--385}
% }
%
% @article{milgram1994,
%   author  = {Milgram, Paul and Kishino, Fumio},
%   title   = {A Taxonomy of Mixed Reality Visual Displays},
%   journal = {IEICE Transactions on Information and Systems},
%   year    = {1994}
% }
%
% @misc{webxr2020,
%   author = {{W3C}},
%   title  = {WebXR Device API},
%   year   = {2020},
%   url    = {https://www.w3.org/TR/webxr/}
% }
%
% @inproceedings{freitas2025benchmark,
%   author = {Freitas, Jorge Luiz Cunha de},
%   title  = {Facilitando a Criação de Experiências Educacionais em
%             Realidade Aumentada por Usuários Não Técnicos: Um Mapeamento
%             Exploratório e Avaliação Técnica de Ferramentas de IA e WebAR},
%   booktitle = {[Nome do Evento/Disciplina]},
%   year = {2025}
% }

% Apêndices (opcional)
\begin{apendicesenv}
  \partapendices

  \chapter{Questionário SUS Aplicado}
  \label{ap:sus}
  % [Inserir o questionário SUS adaptado para o ContAR]

  \chapter{Ficha de Registro de Problemas do Expert Walkthrough}
  \label{ap:ficha}
  % [Inserir a ficha de coleta de dados]

  \chapter{Termo de Consentimento Livre e Esclarecido (TCLE)}
  \label{ap:tcle}
  % [Inserir o TCLE utilizado nas sessões]

\end{apendicesenv}

\end{document}
